% Source: https://artofproblemsolving.com/community/c3h2280289_2000th_post
% Stable identity: aops-topic-2280289-q02
\textbf{Product of all divisors.} If $n$ has $d$ positive divisors, they pair up as $t$ and $\dfrac {n}{t}$, and the product of each pair is $n$. Hence the product of all divisors is $n^{\dfrac {d}{2}}$ (if $n$ is a square, the divisor $\sqrt n$ is paired with itself, and the formula still holds). The product of the proper divisors is therefore $n^{\dfrac {d}{2}-1}$.

\textbf{The condition.} For $n>1$, $n^{\dfrac {d}{2}-1}=n$ holds exactly when $\dfrac {d}{2}-1=1$, that is, $d=4$. The numbers with exactly $4$ divisors are $p^3$ and $pq$ for primes $p\ne q$. A perfect cube of this kind must be $p^3$, and the smallest is $\boxed{8}$.

Check: the proper divisors of $8$ are $1,2,4$, and $1\cdot2\cdot4=8$.

(If ``divisors'' included $n$ itself, the product $n^{\dfrac {d}{2}}$ would equal $n$ only for $d=2$, that is, for primes, and no prime is a perfect cube; only $n=1$ would remain.)
